% (c) 2026 Stephanie Märklin, OST Ostschweizer Fachhochschule
%
% !TEX root = ../../buch.tex
% !TEX encoding = UTF-8
%
% Rolle im Paper: Beleg
%
% Stichworte:
% 17.6 Rechenbeispiel und Vergleich realer Flugzeuge
% - Vorgehen: Systemmatrix mit Zahlenwerten (Nelson Tab. 5.2,
%   Flugzustand nennen); Eigenwerte numerisch berechnen (MATLAB)
% - Feststellung: alle Realteile negativ -> Seitenbewegung stabil
% - Vergleich Business Jet (Stengel \cite): Spiral positiv ->
%   instabil, aber langsam; Dutch Roll schwächer gedämpft

\section{Rechenbeispiel und Vergleich realer Flugzeuge
\label{aircraft:section:6-rechenbeispiel}}
\kopfrechts{Rechenbeispiel}
Wir setzen die Derivative eines einmotorigen Propellerflugzeugs im
Horizontalflug auf Meereshöhe bei 176\,ft/s ein, die Werte stammen
aus Nelson~\cite{aircraft:nelson}, Tabelle~5.2.
Die Systemmatrix lautet damit
\begin{equation}
    A
    =
    \begin{pmatrix}
        -0.254 &  0     & -1.0   & 0.182 \\
        -16.02 & -8.40  &  2.19  & 0     \\
        4.488 & -0.350 & -0.760 & 0     \\
        0     &  1     &  0     & 0
    \end{pmatrix}.
    \label{aircraft:eqn:matrix-ga}
\end{equation}
Die Nullstellen der charakteristischen Gleichung berechnen wir
numerisch mit MATLAB.
% OFFEN: Folie 12 -- Befehl/Skript nennen, falls im Anhang gewünscht
Das ergibt die vier Eigenwerte
\begin{align}
    \lambda_1 &= -8.435, \notag \\
    \lambda_2 &= -0.00877, \notag \\
    \lambda_{3,4} &= -0.487 \pm 2.335\,i.
    \label{aircraft:eqn:eigenwerte-ga}
\end{align}
Tatsächlich stimmen sie mit den von Nelson~\cite{aircraft:nelson}
angegebenen Werten überein.
Die beiden reellen Eigenwerte gehören zur Rollbewegung und zur
Spiralbewegung, das komplexe Paar ist die Dutch Roll.
Alle vier Realteile sind negativ, die Seitenbewegung dieses
Flugzeugs ist also stabil.

Zum Vergleich betrachten wir einen Business Jet, dessen Eigenwerte
Stengel~\cite{aircraft:stengel} angibt:
\begin{align}
    \lambda_1 &= -1.2, \notag \\
    \lambda_2 &= +0.00883, \notag \\
    \lambda_{3,4} &= -0.116 \pm 1.39\,i.
    \label{aircraft:eqn:eigenwerte-jet}
\end{align}
Der Eigenwert der Spiralbewegung ist positiv, die Seitenbewegung
dieses Flugzeugs ist nach Abschnitt~\ref{aircraft:subsection:realteil-imaginaerteil}
also instabil.
Der Realteil der Dutch Roll ist negativ, aber deutlich näher bei
null als beim Propellerflugzeug.

Die Abbildung~\ref{aircraft:fig:komplexe-ebene} zeigt die
Eigenwerte beider Flugzeuge in der komplexen Ebene.
Der Abstand eines Eigenwerts von der imaginären Achse ist seine
Dämpfung, sein Abstand von der reellen Achse seine Frequenz.
Aus beiden folgen die Periode $T = 2\pi/\omega$ der Schwingung und
die Zeit $t_{1/2} = \ln 2 / |\sigma|$, in der die Amplitude auf die
Hälfte fällt.
Die Tabelle~\ref{aircraft:table:dutchroll} enthält diese Werte für
die Dutch Roll beider Flugzeuge.

%\input{papers/aircraft/fig/komplexe-ebene.tex}

\begin{table}
    \centering
    \begin{tabular}{lrrrr}
        \hline
        Flugzeug & $\sigma$ & $\omega$ & $T$ [s] & $t_{1/2}$ [s] \\
        \hline
        Propellerflugzeug & $-0.487$ & $2.335$ & $2.7$ & $1.4$ \\
        Business Jet      & $-0.116$ & $1.39$  & $4.5$ & $6.0$ \\
        \hline
    \end{tabular}
    \caption{Dämpfung, Frequenz, Periode und Halbwertszeit der Dutch
    Roll für die beiden Flugzeuge.
    \label{aircraft:table:dutchroll}}
\end{table}